\documentclass[book.tex]{subfiles}

\begin{document}
\chapter{Quantum Reinforcement Learning}
\label{chap:quantum_rl}
\noindent\rule{11cm}{0.4pt} \\
\textbf{Prerequisites:} \autoref{chap:molecular_hamiltonian}, \autoref{chap:dft}, \autoref{chap:hartree_fock} \\
\textbf{Difficulty Level:} ** \\
\noindent\rule{11cm}{0.4pt}
\vspace{5mm}

"Quantum machine learning has been identified as one of the key fields that could reap advantages from near-term quantum devices, next to optimization and quantum chemistry. Research in this area has focused primarily on variational quantum algorithms, and several proposals to enhance supervised, unsupervised and reinforcement learning algorithms with quantum computing have been put forward. Out of the three, RL is the least studied and it is still an open question whether near-term quantum algorithms can be competitive with state-of-the-art classical approaches based on neural networks even on simple benchmark tasks. In this talk, I will introduce a variational quantum algorithm for deep Q-learning and explain which architectural choices of the quantum model are crucial to make it competitive with its classical counterpart on a benchmark learning task."

References \cite{skolik2021quantum}.
\end{document}
